\section{Sondervorlesung Komplexe Zahlen}
%
%
%
\subsection{Einführung}
$\mathbb{N}_{0}\subset\mathbb{Z}\subset\mathbb{Q}\subset\mathbb{R}$
%
%
%
\subsubsection{$\mathbb{N}_{0}\rightarrow\mathbb{Z}$:}
$\mathbb{Z}:=\mathbb{N}_{0}\times\mathbb{N}_{0}$ (Betrachte Paare von natürlichen Zahlen)\\
\qquad\\
Äquivalenzklassen:\\
$(n,m)\sim(n',m'):\Leftrightarrow n+m' = n'+m \Leftrightarrow n-m=n'-m'$\\
$(n-m)\cdot(n'-m')=nn'+mm'-nm'-n'm$\\
$(7,4)\sim(3,0)$\\
$[n,m]:=\{(n',m')|(n',m')\sim(n,m)\}$\\
\qquad\\
Definiere, was auf diese Menge "`+"' und "`$\cdot$"' sein soll:\\
$[n,m]+[n',m']:=[n'+m,n+m']$\\
$[n,m]\cdot[n',m']:=[nn'+mm',nm'+n'm]$\\
$n:=[n,0]$\qquad $-n:=[0,n]$
%
%
%
\subsubsection{$\mathbb{Z}\rightarrow\mathbb{Q}$:}
$\mathbb{Q}=\mathbb{Z}\times\mathbb{Z}\diagdown N$\\
$(a,b)\sim(a',b'):\Leftrightarrow ab'=ba' \qquad \frac{a}{b}=\frac{a'}{b'}\Leftrightarrow ab'=a'b$\\
$\frac{a}{b}:=[a,b]$\\
Addition: $[a,b]+[c,d]:=[ad+bc,bd]$\\
Multiplikation: $[a,b]+[c,d]:=[ac,bd]$\\
$\Rightarrow$ Multiplikation in $\mathbb{Q}$ ist genauso definiert wie Addition in $\mathbb{Z}$!
%
%
%
\subsubsection{$\mathbb{Q}\rightarrow\mathbb{R}$:}
$\mathbb{R}:=\dotsc$\\
Konstruktion sehr kompliziert, über Zahlenfolgen in $\mathbb{Q}$.
$x^{2}+1=0$ hat in $\mathbb{R}$ keine Lösung\\
$x^{2}\geq0 \, \forall x\in \mathbb{R}\Rightarrow x^{2}+1\geq1>0$\\
Cauchi-Folge: Abstand zwischen Folgeglieder wird immer kleiner trotzdem u.U. kein Grenzwert in $\mathbb{Q}$
%
%
%
\subsubsection{$\mathbb{R}\rightarrow\mathbb{C}$:}
$x^{2}\mp 0 \, \forall x\in\mathbb{R}\Rightarrow x^{2}+\geq>0$\\
Wir suchen einen Körper $\mathbb{C}$ mit 
\begin{enumerate}
	\item $\mathbb{R}\subset\mathbb{C}$
	\item Die Gleichung hat (mindestens) eine Lösung.
\end{enumerate}
Als Menge ist $\mathbb{C}:=\mathbb{R}\times\mathbb{R}=\{(x,y)|x,y\in\mathbb{R}\}$\\
Wir definiere eine Addition durch folgende Vorschrift:\\
$(a,b)+(c,d):=(a+c,b+d)$\\
und eine Multiplikation durch:\\
$(a,b)(c,d):=(ac-bd,ad+bc)$
%
%
%
\subsubsection{Behauptung:}
$(\mathbb{C},+,\cdot)$ ist eine Körper mit $0=(0,0)$ (Nullelement)\\
$1\mathop{=}\limits^{\text{*}}(1,0)$\\
\qquad\\
*($(a,b)(x,y)=(a,b)$\\
$=(ax-by,ay+bx)$\\
$ax-by=a$\\
$bx+ay=b \leadsto(x,y)=(1,0)$)\\
\qquad\\
Betrachtet man Elemente der Form $(x,0) \in \mathbb{C}$, dann gilt:\\
$(a,0)+(b,0)=(a+b,0)$\\
$(a,0)\cdot(c,0)=(ac,0)$\\
\qquad\\
Wir können also $\mathbb{R}$ als Unterkörper von $\mathbb{C}$ (bzw. $\mathbb{C}$ als Körpererweiterung von $\mathbb{R}$) auffassen durch die Identifizierung:\\
$\mathbb{R}\rightarrow\mathbb{C}$\\
$a\rightarrow(a,0)$\\
Was ist Nullstelle von $x^{2}+1=0$?
 \begin{align}
	x&=(a,b)\\
    (0,0)=0=x^{2}+1 &=& (a,b)^{2}+(1,0) \\
	&= (a,b)(a,b)+(1,0) \leadsto \qquad \text{i.) }0\mathop{=}\limits^{\text{!}}a^{2}-b^{2}+1\\
	&= (a^{2}-b^{2},2ab)+(1,0) \leadsto \quad \text{ii.) } 0 \mathop{=}\limits^{\text{!}} 2ab\\
	&= (a^{2}-b^{2}+1,2ab) \leadsto \qquad \quad \text{iii.) } a=0 \text{ oder } b=0\\
  \end{align}
\begin{itemize}
	\item $b=0 \Rightarrow 0=a^{2}+1 \quad \lightning$
	\item $a=0 \Rightarrow 0=-b^{2}+1\Rightarrow b = \pm\sqrt{1} = 1$
\end{itemize}
$
x=\left\{
\begin{array}{l}     
    (0,1)\\
    (0,-1)
\end{array}\right.
$\\
\qquad\\
Wir schreiben ab jetzt:\\
$i:=(0,1)$\\
$-i:=(0,-1)$\\
und wir schreiben: $a:=(a,0)$ für $a\in\mathbb{R}$\\
dann lässt sich jede komplexe Zahl wie folgt schreiben:\\
$z=(a,b)\in\mathbb{C}=\underline{a+bi}$\\
.[ 
\begin{align}
a+bi &= (a,0)+(b,0)\cdot(0,1)\\
&= (a,0)+(0,b\cdot1)\\
&= (a,b)\\
\end{align}
Eine komplexe Zahl ist die Information, zweier reellen Zahlen $a$ und $b$ zu haben und ein Element $i$, das mit sich selbst multipliziert $-1$ ergibt.\quad ]\\
$\mathbb{C}$ ist ein 2-dimensionaler $\mathbb{R}$-Vektorraum mit der Basis $(1,i)$.\\
$p(x)=0, \, p\in \mathbb{C}[x] \rightarrow$ Fundamentalsatz der Algebra: "`Jedes Polynom besitzt mindestens eine Nullstellen. Jedes Polynom hat so viele Nullstellen, wie sein Grad ist. "'\\
\qquad\\
\includegraphics[width=0.5\textwidth]{mainmatter/chapter3/pics/komplexkoord.png}
%
%
%
\subsection{Definitionen:}
Für die komplexe Zahl $z=x+iy \forall x,y \in \mathbb{R}$ definieren wir:
\begin{enumerate}
\item $RE(z):=x$ Realteil von $z$\\
	$Im(z):=y$ Imaginärteil von $z \rightarrow$ (Nur $y$, nicht $i$, wichtig!)\\
	$\leadsto z=Re(z)+i\,Im(z)$
\item $\overline{z}:=x-iy=Re(z)-i\,Im(z)$ zu $z$ komplex konjugierte Zahl. \\
	.[ $z\cdot\overline{z} = (x+iy)(x-iy) = x^{2}+iyx-ixy-i^{2}y^{2} = x^{2}+y^{2}\geq0$ .]
\item $\vert z\vert :=\sqrt{z\overline{z}}$ Betrag von $z$, $\sqrt{}$ = positive reelle Wurzel.\\
	\begin{enumerate}
	\item $\overline{z}=z$
	\item $\overline{z+w}=\overline{z}+\overline{w}$  (Konjugation ist lineare Abbildung 
	$\mathbb{R}^{2}\rightarrow\mathbb{R}^{2}$)\\
	$\overline{z\cdot w} = \overline{z}\cdot\overline{w}$
	\item $z=z\Leftrightarrow Im(z)\Leftrightarrow z\in \mathbb{R}$\\
	$z=-\overline{z}\Leftrightarrow Re(z)=0\Leftrightarrow z \in i\cdot\mathbb{R}$
	\item $Re(z)=\frac{1}{2}(z+\overline{z}), \,Im(z)=\frac{1}{2i}(z-\overline{z})$
	\item$\vert z\vert \geq 0, \,  \vert z \vert = 0 \Leftrightarrow z = 0$
	\item $\vert z+w \vert \leq \vert z \vert + \vert w \vert$ (Dreiecksungleichung)
	\end{enumerate}
\end{enumerate}
%
%
%
\subsubsection{Inverses von $z$:}
$z^{-1}=\frac{1}{z}=\frac{\overline{z}}{\vert z \vert^{2}}$\\
Beispiel:\\
$(1-2)^{-1}=\frac{1+i}{2}$\\
$\begin{pmatrix} 1 \\ -1 \end{pmatrix}\Rightarrow\vert\sqrt{2}\vert$\\
$\vert\sqrt{2}\vert^{2}=2$\\
Vorzeichen vo dem i muss umgekehrt werden!
\begin{figure}
\centering
\includegraphics[width=0.35\textwidth]{mainmatter/chapter3/pics/zrechnungen.png}
\caption{$z$ und $\overline{z}$ gespiegelt an der x-Achse}
\end{figure}
\begin{align}
z &= x+iy\\
&= r\cos\varphi + ir\sin\varphi \rightarrow \text{Polarkoordinatenschreibweise}\\
&= \vert z \vert \mathop{\underbrace{(\cos\varphi+i\sin\varphi)}}\limits_{\text{Alle Zahlen mit } \vert 1 \vert \text{ (liegen auf dem Einheitskreis) }}\\
&\qquad \qquad \qquad \mathop{\Downarrow}\limits_{exp(i\varphi)=e^{i\varphi}}\\
\end{align}
$\Rightarrow$ Realteil bleibt gleicht, Imaginärteil wird ins negative verkehrt.\\
komplexe Zahlen sind nicht sortierbar:\\
$a>0\Rightarrow a^{2}>0$\\
\begin{align}
i^{2} &= -1 < 0\\
(-i)^{2} &= -1<0\\
i &=\sqrt{-1} \Rightarrow \text{stimmt, aber wir kennen Hintergrund noch nicht } \Rightarrow \text{ nicht so schreiben! }\\
i\cdot i &= \sqrt{-1}\cdot\sqrt{-1} \neq \sqrt{(-1)\cdot(-1)}=\sqrt{1}=1
\end{align}